% =====================================================================
% Informe: Implementación de la arquitectura IoT del Gemelo Digital L418
% Compilar en Overleaf con pdfLaTeX
% =====================================================================
\documentclass[11pt,a4paper]{article}

\usepackage[utf8]{inputenc}
\usepackage[T1]{fontenc}
\usepackage[spanish,es-noshorthands,es-tabla]{babel}
\usepackage{lmodern}
\usepackage[margin=2.5cm]{geometry}
\usepackage{amsmath,amssymb}
\usepackage{booktabs,tabularx,longtable,array}
\usepackage{xcolor}
\usepackage{listings}
\usepackage[most]{tcolorbox}
\usepackage{tikz}
\usetikzlibrary{arrows.meta,positioning,fit,backgrounds,calc}
\usepackage{enumitem}
\usepackage{float}
\usepackage{caption}
\usepackage{fancyhdr}
\usepackage[hidelinks]{hyperref}

% ---------- Estilos ----------
\definecolor{fondocod}{gray}{0.96}
\definecolor{azul}{RGB}{31,78,121}
\definecolor{naranja}{RGB}{200,100,20}

\lstset{
  basicstyle=\ttfamily\small,
  backgroundcolor=\color{fondocod},
  frame=single, rulecolor=\color{gray!50},
  breaklines=true, breakatwhitespace=false,
  columns=fullflexible, keepspaces=true,
  showstringspaces=false, tabsize=2,
  xleftmargin=4pt, xrightmargin=4pt,
  commentstyle=\color{gray}\itshape,
  keywordstyle=\color{azul}\bfseries
}
\lstdefinestyle{bash}{language=bash}
\lstdefinestyle{cpp}{language=C++}
\lstdefinestyle{plano}{language={}}

\newtcolorbox{completar}{
  colback=naranja!6, colframe=naranja!80!black,
  fonttitle=\bfseries, title=Completar antes de usar,
  left=6pt, right=6pt, top=4pt, bottom=4pt
}
\newtcolorbox{nota}[1][Nota]{
  colback=azul!5, colframe=azul!70!black,
  fonttitle=\bfseries, title=#1,
  left=6pt, right=6pt, top=4pt, bottom=4pt
}

\newcommand{\arch}[1]{\texttt{#1}}

\pagestyle{fancy}
\fancyhf{}
\fancyhead[L]{\small Gemelo Digital -- Sal\'on L418}
\fancyhead[R]{\small UTEC}
\fancyfoot[C]{\thepage}

\setlength{\parskip}{4pt}
\setlength{\emergencystretch}{3em}   % evita líneas desbordadas con nombres de variables largos
\setlist{itemsep=2pt, topsep=3pt}

% =====================================================================
\begin{document}

\begin{titlepage}
  \centering
  {\large Universidad de Ingenier\'ia y Tecnolog\'ia -- UTEC\par}
  \vspace{2.5cm}
  {\Huge\bfseries Arquitectura IoT para el\\[4pt] Gemelo Digital del Sal\'on L418\par}
  \vspace{0.8cm}
  {\Large Implementaci\'on en servidor local y migraci\'on a la nube\par}
  \vspace{2.5cm}
  {\large \textbf{Integrantes:}\par}
  \vspace{0.2cm}
  {\large Abel Romulo Chuquillanqui Marav\'i\\ Sebastian Falvy Mendoza\\ Juan Diego Mej\'ia Armasa\\ Deyver Ricardo Fuertes Palpa\par}
  \vspace{0.6cm}
  {\large \textbf{Profesor:} Prof.\ Figueroa\par}
  \vspace{0.3cm}
  {\large \textbf{Curso:} PI3103 -- Proyectos Interdisciplinarios 3\par}
  \vfill
  {\large Lima, noviembre de 2026\par}
\end{titlepage}

\begin{abstract}
Este informe documenta, paso a paso, la implementaci\'on de una arquitectura de
Internet de las Cosas (IoT) que alimenta el gemelo digital del sal\'on L418. Siete nodos
basados en ESP32 miden la concentraci\'on de CO$_2$ (MH-Z1911A), el nivel sonoro (INMP441),
gases combustibles y humo (MQ-2), la iluminaci\'on (BH1750), la ocupaci\'on mediante una
c\'amara t\'ermica (MLX90640), el estado de la puerta (KY-024) y la temperatura, humedad y
presi\'on (BME280). Cada nodo procesa sus mediciones localmente, detecta anomal\'ias y
transmite por ESP-NOW a un ESP32 central (gateway), que publica los datos por MQTT hacia un
servidor con Mosquitto, Telegraf, InfluxDB y Grafana. Los datos se visualizan en tiempo real,
se exportan a CSV para cualquier rango de tiempo y alimentan el gemelo digital 2D del
sal\'on. El sistema se valid\'o de punta a punta con un servidor local, cada nodo se verific\'o
en simulaci\'on, y se prepar\'o la versi\'on de producci\'on para la nube, con HTTPS, MQTT cifrado
y acceso p\'ublico de solo lectura. Los archivos citados se encuentran en el repositorio
indicado en los anexos.
\end{abstract}

\tableofcontents
\newpage

% =====================================================================
\section{Objetivo y alcance}

El objetivo es que la informaci\'on de los sensores del sal\'on L418 llegue a un
dashboard en tiempo real, accesible para la comunidad universitaria, desde el cual el
usuario pueda exportar los datos del rango de tiempo que desee analizar. Esta
infraestructura de datos es la base del gemelo digital.

El alcance de este informe cubre:
\begin{enumerate}
  \item Los siete nodos sensores del sal\'on y su procesamiento en el borde.
  \item El enlace inal\'ambrico ESP-NOW y el gateway ESP32.
  \item El servidor local en una laptop (validaci\'on completa) y el dashboard final.
  \item La migraci\'on a la nube (Azure for Students) y la versi\'on de producci\'on del servidor.
  \item Los requisitos para alimentar el gemelo digital 2D del sal\'on.
\end{enumerate}

Los siete sensores pertenecen al sal\'on: hay una sola medici\'on de cada variable y todos
los nodos se comunican con un \'unico gateway. Los datos son compartidos por todo el
proyecto; cada integrante aporta desde su responsabilidad.

\subsection{Equipo y responsabilidades}

\begin{table}[H]
\centering
\caption{Responsabilidades del equipo.}
\begin{tabularx}{\textwidth}{@{}lX@{}}
\toprule
\textbf{Integrante} & \textbf{Responsabilidad} \\
\midrule
Abel R.\ Chuquillanqui Marav\'i & Integraci\'on de hardware (soldadura y conexionado de los sensores) y arquitectura IoT: nodos de CO$_2$ y ruido, comunicaci\'on ESP-NOW/MQTT, gateway, servidor y dashboard \\
Sebastian Falvy Mendoza & Gemelo digital 2D (aplicaci\'on web Web-PI3) \\
Juan Diego Mej\'ia Armasa & Gemelo digital 2D (aplicaci\'on web Web-PI3) \\
Deyver R.\ Fuertes Palpa & Dise\~no de las carcasas de los m\'odulos y gesti\'on del proyecto (Asana) \\
\bottomrule
\end{tabularx}
\end{table}

% =====================================================================
\section{Arquitectura general}

La arquitectura se organiza en cuatro capas. Cada capa tiene una sola responsabilidad, lo
que permite cambiar una (por ejemplo, el servidor) sin modificar las dem\'as.

\begin{figure}[H]
\centering
\begin{tikzpicture}[
  font=\small,
  caja/.style={draw, rounded corners=3pt, align=center, minimum height=1.0cm,
               minimum width=3.0cm, fill=#1},
  caja/.default=white,
  flecha/.style={-{Stealth[length=2.2mm]}, thick, gray!70!black},
  etiq/.style={font=\scriptsize\itshape, fill=white, inner sep=1pt}
]
  % Capa 1: sensores
  \node[caja=teal!12] (co2)   {Nodos 1, 2 y 3\\CO$_2$, ruido, gas};
  \node[caja=teal!12, right=0.6cm of co2] (aud) {Nodos 4 y 7\\luz, ambiente};
  \node[caja=teal!12, right=0.6cm of aud] (otr) {Nodos 5 y 6\\personas, puerta};

  % Capa 2: gateway
  \node[caja=orange!15, below=1.3cm of aud, minimum width=4.2cm] (gw)
       {Gateway ESP32 v4\\ESP-NOW $\rightarrow$ WiFi};

  \draw[flecha] (co2.south) |- ($(gw.north)+(-1.2,0.55)$) -- ($(gw.north)+(-1.2,0)$);
  \draw[flecha] (aud.south) -- (gw.north);
  \draw[flecha] (otr.south) |- ($(gw.north)+(1.2,0.55)$) -- ($(gw.north)+(1.2,0)$);
  \node[etiq] at ($(aud.south)!0.45!(gw.north)+(1.05,0)$) {ESP-NOW};

  % Capa 3: servidor
  \node[caja=violet!10, below=1.6cm of gw, xshift=-4.35cm, minimum width=2.4cm] (mq)
       {Mosquitto\\broker MQTT};
  \node[caja=violet!10, right=0.5cm of mq, minimum width=2.4cm] (tg) {Telegraf\\traductor};
  \node[caja=violet!10, right=0.5cm of tg, minimum width=2.4cm] (ix) {InfluxDB\\series de tiempo};
  \node[caja=violet!10, right=0.5cm of ix, minimum width=2.4cm] (gf) {Grafana\\dashboard};
  \draw[flecha] (mq) -- (tg);
  \draw[flecha] (tg) -- (ix);
  \draw[flecha] (ix) -- (gf);

  \begin{scope}[on background layer]
    \node[draw, dashed, rounded corners=5pt, inner sep=8pt, fit=(mq)(tg)(ix)(gf),
          label={[font=\scriptsize\itshape]above right:Servidor (laptop o nube) -- Docker}] (srv) {};
  \end{scope}

  \draw[flecha] (gw.south) |- ($(mq.north)+(0,0.75)$) -- (mq.north);
  \node[etiq] at ($(gw.south)+(0,-0.45)$) {MQTT por WiFi};

  % Capa 4: usuarios
  \node[caja=gray!12, below=1.0cm of srv] (csv) {Exportar CSV\\rango de tiempo};
  \node[caja=gray!12, left=0.5cm of csv] (dash) {Dashboard p\'ublico\\(navegador)};
  \node[caja=gray!12, right=0.5cm of csv] (gem) {Gemelo digital 2D\\(Web-PI3)};
  \draw[flecha] (srv.south) -- ++(0,-0.45) -| (dash.north);
  \draw[flecha] (srv.south) -- (csv.north);
  \draw[flecha] (srv.south) -- ++(0,-0.45) -| (gem.north);
\end{tikzpicture}
\caption{Arquitectura por capas del sistema.}
\label{fig:arquitectura}
\end{figure}

\begin{table}[H]
\centering
\caption{Funci\'on de cada capa.}
\begin{tabularx}{\textwidth}{@{}lX@{}}
\toprule
\textbf{Capa} & \textbf{Funci\'on} \\
\midrule
Nodos sensores & Cada ESP32 lee su sensor, calcula indicadores y detecta anomal\'ias. Env\'ia solo
indicadores, nunca datos crudos (por ejemplo, el nodo de audio no transmite el sonido). \\
Enlace ESP-NOW & Protocolo inal\'ambrico directo entre ESP32. No requiere router ni red WiFi. \\
Gateway & \'Unico dispositivo conectado a WiFi. Asigna la hora real (NTP), convierte los
paquetes a JSON y los publica por MQTT. \\
Servidor & Mosquitto recibe los mensajes MQTT, Telegraf los traduce, InfluxDB los almacena
como series de tiempo y Grafana los muestra y permite exportarlos. \\
\bottomrule
\end{tabularx}
\end{table}

\begin{table}[H]
\centering
\caption{Nodos del sal\'on L418.}
\small
\begin{tabularx}{\textwidth}{@{}clllX@{}}
\toprule
\textbf{Nodo} & \textbf{Sensor} & \textbf{Interfaz} & \textbf{Env\'io} & \textbf{Ubicaci\'on} \\
\midrule
1 & MH-Z1911A (CO$_2$) & UART & 5 s & Techo, centro \\
2 & INMP441 (ruido) & I$^2$S & 1 s & Techo, centro \\
3 & MQ-2 (gas y humo) & Anal\'ogica & 2 s & Techo, separado de los dem\'as \\
4 & BH1750 (iluminaci\'on) & I$^2$C & 10 s + evento & Techo, centro \\
5 & MLX90640 (personas) & I$^2$C & 5 s + evento & Marco superior de la puerta (2.3 m) \\
6 & KY-024 (puerta) & Anal\'ogica & 10 s + evento & Marco de la puerta, im\'an en la hoja \\
7 & BME280 (T, HR, P) & I$^2$C & 10 s & Techo, centro \\
\bottomrule
\end{tabularx}
\end{table}

\subsection{Decisiones de dise\~no}

\begin{itemize}
  \item \textbf{ESP32 como gateway y no como dashboard.} Un ESP32 tiene unos 520~KB de RAM:
        no puede almacenar semanas de datos de varios sensores ni atender a muchos usuarios.
        Por eso reenv\'ia los datos y el dashboard vive en un servidor.
  \item \textbf{MQTT en lugar de archivos CSV.} El CSV es un formato de exportaci\'on, no de
        transporte. Enviar archivos impide el tiempo real y obliga a fusionarlos a mano. Con
        MQTT los datos fluyen continuamente y el CSV se genera al final, para el rango que
        elija el usuario.
  \item \textbf{ESP-NOW en lugar de Bluetooth.} El ESP32 tiene una sola radio compartida
        entre WiFi y Bluetooth. ESP-NOW usa la misma radio que WiFi, tiene mayor alcance y
        es m\'as simple para varios emisores hacia un receptor.
  \item \textbf{Un solo gateway para el sal\'on.} Los siete nodos usan ESP-NOW hacia el
        mismo gateway. As\'i, TI solo debe autorizar un dispositivo en la red WiFi y
        existe un \'unico usuario MQTT.
  \item \textbf{Procesamiento en el borde.} Enviar audio crudo requerir\'ia
  \begin{equation}
    16\,000~\tfrac{\text{muestras}}{\text{s}} \times 16~\tfrac{\text{bits}}{\text{muestra}}
    = 256~\text{kbps},
  \end{equation}
  inviable para un enlace inal\'ambrico de bajo consumo. El nodo env\'ia en cambio un
  paquete de 38~bytes por segundo.
\end{itemize}

% =====================================================================
\section{Materiales y herramientas}

\begin{table}[H]
\centering
\caption{Hardware utilizado.}
\begin{tabularx}{\textwidth}{@{}lX@{}}
\toprule
\textbf{Componente} & \textbf{Uso} \\
\midrule
ESP32 DevKit V1 + MH-Z1911A & Nodo 1: CO$_2$ (UART) \\
ESP32 DevKit V1 + INMP441 & Nodo 2: nivel sonoro (I$^2$S) \\
ESP32 DevKit V1 + MQ-2 & Nodo 3: gases combustibles y humo (anal\'ogico, 5 V) \\
ESP32 DevKit V1 + GY-302 BH1750 & Nodo 4: iluminaci\'on (I$^2$C) \\
ESP32 DevKit V1 + MLX90640 (55$^\circ$) & Nodo 5: conteo de personas (I$^2$C) \\
ESP32 DevKit V1 + KY-024 + im\'an & Nodo 6: estado de la puerta (anal\'ogico) \\
ESP32 DevKit V1 + GY-BME280 & Nodo 7: temperatura, humedad y presi\'on (I$^2$C) \\
ESP32 DevKit V4 & Gateway (nodo central) \\
Resistencias 10 k$\Omega$ y 15 k$\Omega$ (2 de cada una) & Divisores de tensi\'on del MQ-2 \\
Laptop con Windows 11 & Servidor local y programaci\'on \\
\bottomrule
\end{tabularx}
\end{table}

\begin{table}[H]
\centering
\caption{Software utilizado.}
\begin{tabularx}{\textwidth}{@{}lX@{}}
\toprule
\textbf{Software} & \textbf{Uso} \\
\midrule
Arduino IDE + n\'ucleo ESP32 & Programaci\'on de los tres ESP32 \\
Librer\'ia PubSubClient (Nick O'Leary) & Cliente MQTT en el gateway \\
Librer\'ias Adafruit MLX90640, Adafruit BME280 y BH1750 (C.\ Laws) & Lectura de los sensores I$^2$C \\
Docker Desktop + WSL~2 & Ejecuci\'on de los servicios del servidor \\
\texttt{eclipse-mosquitto:2} & Broker MQTT \\
\texttt{telegraf:1.32} & Ingesta MQTT $\rightarrow$ InfluxDB \\
\texttt{influxdb:2.7} & Base de datos de series de tiempo \\
\texttt{grafana/grafana:11.3.0} & Dashboard y exportaci\'on CSV \\
\texttt{caddy:2.8} & Proxy HTTPS (solo en la nube) \\
Microsoft Azure for Students & M\'aquina virtual en la nube \\
\bottomrule
\end{tabularx}
\end{table}

% =====================================================================
\section{Fase 1: Nodos sensores}

\subsection{Formato com\'un de datos}

Los tres ESP32 comparten el archivo \arch{paquetes.h}, que define la estructura binaria de
cada paquete. Se eligi\'o un formato binario porque ESP-NOW admite como m\'aximo 250~bytes
por paquete y porque el gateway no necesita interpretar texto. Las anomal\'ias viajan como
bits dentro de un byte.

\begin{completar}
\arch{paquetes.h}: no se modifica. Debe ser \textbf{id\'entico} en las carpetas de los siete
nodos, \arch{gateway} y \arch{gateway\_nube}. Si se cambia, hay que actualizar todas las copias y
volver a programar al menos el gateway y los nodos cuya estructura cambi\'o.
\end{completar}

\begin{table}[H]
\centering
\caption{Cabecera com\'un de todos los paquetes (11~bytes).}
\begin{tabular}{@{}lll@{}}
\toprule
\textbf{Campo} & \textbf{Tipo} & \textbf{Descripci\'on} \\
\midrule
\texttt{version} & \texttt{uint8} & Versi\'on del protocolo (1) \\
\texttt{tipo} & \texttt{uint8} & 1 = CO$_2$, 2 = audio \\
\texttt{nodo\_id} & \texttt{uint8} & Identificador \'unico del nodo \\
\texttt{seq} & \texttt{uint32} & Contador para detectar paquetes perdidos \\
\texttt{uptime\_ms} & \texttt{uint32} & Tiempo desde el encendido del nodo \\
\bottomrule
\end{tabular}
\end{table}

\begin{table}[H]
\centering
\caption{Tipos de paquete y anomal\'ias (bit 0, 1, 2, 3).}
\small
\begin{tabularx}{\textwidth}{@{}clX@{}}
\toprule
\textbf{Tipo} & \textbf{Nombre} & \textbf{Anomal\'ias por bit} \\
\midrule
1 & \texttt{co2} & \texttt{FALLO\_SENSOR}, \texttt{CO2\_PICO\_RAPIDO}, \texttt{CO2\_ALERTA\_VENTILACION}, \texttt{CO2\_CRITICO} \\
2 & \texttt{audio} & \texttt{FALLO\_AUDIO}, \texttt{RUIDO\_IMPACTO}, \texttt{RUIDO\_FUERTE\_SOSTENIDO}, \texttt{RUIDO\_MECANICO\_CONSTANTE} \\
3 & \texttt{gas} & \texttt{FALLO\_SENSOR}, \texttt{GAS\_DETECTADO}, \texttt{GAS\_CRITICO} \\
4 & \texttt{luz} & \texttt{FALLO\_SENSOR}, \texttt{ILUMINACION\_BAJA} \\
5 & \texttt{termico} & \texttt{FALLO\_SENSOR}, \texttt{CONTEO\_CORREGIDO}, \texttt{OBJETO\_ESTATICO} \\
6 & \texttt{magnetico} & \texttt{FALLO\_SENSOR}, \texttt{PUERTA\_ABIERTA\_PROLONGADA} \\
7 & \texttt{ambiente} & \texttt{FALLO\_SENSOR}, \texttt{TEMP\_FUERA\_CONFORT}, \texttt{HUMEDAD\_ALTA}, \texttt{CAMBIO\_RAPIDO\_TEMP} \\
\bottomrule
\end{tabularx}
\end{table}

\subsection{Nodo de CO$_2$ (MH-Z1911A)}

\textbf{Conexi\'on.} El sensor se comunica por UART2 a 9600~baudios: TX del sensor al
GPIO~16 (RX2) y RX del sensor al GPIO~17 (TX2).

\textbf{Funcionamiento.} Cada 5~s el nodo env\'ia el comando de lectura \texttt{0x86},
valida el encabezado y el checksum de la respuesta, y obtiene la concentraci\'on y la
temperatura interna aproximada:
\begin{equation}
  C_{\text{CO}_2} = 256\,r_2 + r_3 \quad [\text{ppm}],
  \qquad T \approx r_4 - 40 \quad [^\circ\text{C}].
\end{equation}
El periodo de 5~s es adecuado porque el sensor responde en 1 a 2~minutos ($T_{90}$).
Durante los primeros 3~minutos (calentamiento) no se eval\'uan anomal\'ias.

\textbf{Niveles y anomal\'ias.} Los umbrales son 1000~ppm (alerta de ventilaci\'on) y
1500~ppm (cr\'itico), con hist\'eresis de 100~ppm y confirmaci\'on de 3 lecturas
consecutivas para evitar oscilaciones. Adem\'as se calcula la variaci\'on en 30~s
($6 \times 5\,\text{s}$) para detectar subidas r\'apidas mayores a 300~ppm.

\textbf{Archivos:} \arch{nodo\_co2.ino}, \arch{paquetes.h}, \arch{espnow\_nodo.h}
(los tres en la carpeta \arch{nodo\_co2}).

\begin{completar}
En \arch{nodo\_co2.ino}:
\begin{itemize}
  \item \texttt{MAC\_GATEWAY}: MAC del ESP32 v4. En este proyecto:\\
        \texttt{\{0x58, 0x2A, 0xBD, 0x72, 0x03, 0xDC\}} (ver secci\'on~\ref{sec:mac}).
  \item \texttt{SSID\_ROUTER}: nombre exacto de la red WiFi a la que se conecta el gateway.
        Solo se usa para encontrar el canal; el nodo no se conecta a esa red.
  \item \texttt{NODO\_ID}: \texttt{1}. Debe ser \'unico en todo el proyecto.
\end{itemize}
\end{completar}

\subsection{Nodo de audio (INMP441)}

\textbf{Conexi\'on I$^2$S.} SCK al GPIO~26, WS al GPIO~25, SD al GPIO~33; L/R a GND (canal
izquierdo).

\textbf{Procesamiento.} El nodo muestrea a $f_s = 16$~kHz en bloques de $N = 512$ muestras
(32~ms). Para cada bloque calcula el valor RMS y el pico; cada segundo agrega los bloques y
env\'ia los indicadores. El nivel equivalente se calcula como promedio de energ\'ia, no
de decibeles:
\begin{equation}
  L_{eq,T} = 10\log_{10}\!\left(\frac{1}{n}\sum_{i=1}^{n} 10^{L_i/10}\right).
\end{equation}
Por ejemplo, el promedio de 60~dB y 80~dB es aproximadamente 77~dB y no 70~dB, porque el
sonido fuerte domina la energ\'ia.

\textbf{Conversi\'on a dB SPL.} El micr\'ofono entrega niveles en dBFS (relativos al fondo de
escala digital). Con la sensibilidad del INMP441 ($-26$~dBFS a 94~dB SPL):
\begin{equation}
  L_{\text{SPL}} \approx L_{\text{dBFS}} - (-26) + 94 = L_{\text{dBFS}} + 120.
\end{equation}

\textbf{An\'alisis espectral.} Sobre el bloque m\'as fuerte de cada segundo se aplica una
ventana de Hann y una FFT de 512 puntos, con resoluci\'on
\begin{equation}
  \Delta f = \frac{f_s}{N} = \frac{16\,000}{512} = 31.25~\text{Hz}.
\end{equation}
Se obtiene la frecuencia dominante y un \'indice de tonalidad (energ\'ia concentrada en el
pico); un tono persistente durante 30~s se reporta como ruido mec\'anico constante.

\begin{table}[H]
\centering
\caption{Indicadores enviados por el nodo de audio cada segundo.}
\begin{tabular}{@{}ll@{}}
\toprule
\textbf{Campo} & \textbf{Significado} \\
\midrule
\texttt{leq\_db} & Nivel equivalente del segundo (dB SPL estimado) \\
\texttt{lmax\_db} & M\'aximo RMS de bloques de 32~ms \\
\texttt{pico\_db} & Pico de muestra \\
\texttt{base\_db} & Ruido de fondo estimado \\
\texttt{freq\_dom\_hz} & Frecuencia dominante \\
\texttt{tonal} & \'Indice de tonalidad (0 a 1) \\
\texttt{bloques\_fuertes} & Bloques sobre el umbral en el segundo \\
\bottomrule
\end{tabular}
\end{table}

\textbf{Archivos:} \arch{nodo\_audio.ino}, \arch{paquetes.h}, \arch{espnow\_nodo.h}
(carpeta \arch{nodo\_audio}).

\begin{completar}
En \arch{nodo\_audio.ino}:
\begin{itemize}
  \item \texttt{MAC\_GATEWAY}: la misma MAC del gateway.
  \item \texttt{SSID\_ROUTER}: el mismo nombre de red que en el nodo de CO$_2$.
  \item \texttt{NODO\_ID}: \texttt{2}.
  \item \texttt{OFFSET\_SPL\_DB}: valor inicial 120. Ajustarlo comparando con un son\'ometro:
        si el nodo marca $L_{\text{nodo}}$ y el son\'ometro $L_{\text{ref}}$, el nuevo valor es
        $120 - (L_{\text{nodo}} - L_{\text{ref}})$.
\end{itemize}
\end{completar}


\subsection{Nodo de gas y humo (MQ-2)}

\textbf{Conexi\'on.} El MQ-2 se alimenta a 5~V desde el pin VIN del ESP32, porque su calefactor
(33~$\Omega$, consumo menor a 800~mW) necesita ese voltaje para alcanzar la temperatura de
trabajo. Sus salidas AO y DO pueden llegar a 5~V, por lo que cada una pasa por un divisor
($R_1 = 10$~k$\Omega$, $R_2 = 15$~k$\Omega$) antes de entrar al ESP32: AO al GPIO~35 y DO al
GPIO~32, ambos del ADC1.
\begin{equation}
  V_{\text{ESP32}} = V_{\text{m\'odulo}}\,\frac{R_2}{R_1 + R_2} = 5~\text{V}\cdot\frac{15}{25} = 3.0~\text{V}.
\end{equation}

\textbf{Funcionamiento.} El MQ-2 es un sensor semiconductor de SnO$_2$ cuya resistencia baja
ante gases reductores (GLP, humo, alcohol, metano, hidr\'ogeno), sin distinguirlos. A partir
del voltaje de salida se calcula su resistencia y se compara con su valor en aire limpio,
$R_0$, calibrado tras el precalentamiento y guardado en la memoria flash:
\begin{equation}
  R_s = R_L\,\frac{V_C - V_{\text{out}}}{V_{\text{out}}},
  \qquad
  \text{rs\_r0} = \frac{R_s}{R_0}.
\end{equation}
Adem\'as se reporta una estimaci\'on referencial de GLP equivalente con la curva del datasheet,
considerando que el fabricante define su $R_0$ con $R_s/R_0 = 9.83$ en aire limpio:
\begin{equation}
  C_{\text{GLP}} \approx 613.9\,\big(9.83\cdot\text{rs\_r0}\big)^{-2.074}~\text{ppm}.
\end{equation}

\textbf{Niveles.} ALERTA con \texttt{rs\_r0} $< 0.60$ (sale con $> 0.70$) y CR\'ITICO con
\texttt{rs\_r0} $< 0.17$ (sale con $> 0.20$), cada cambio confirmado con 3 lecturas seguidas.
El nivel cr\'itico corresponde aproximadamente a 200~ppm de GLP, el inicio del rango del datasheet.

\textbf{Variables:} \texttt{v\_ao}, \texttt{rs\_r0}, \texttt{ppm\_glp\_est}, \texttt{nivel},
\texttt{do\_alarma}, \texttt{calentando}. \textbf{Archivo:} \arch{nodo\_gas.ino}.

\begin{completar}
En \arch{nodo\_gas.ino}: \texttt{MAC\_GATEWAY}, \texttt{SSID\_ROUTER}, \texttt{NODO\_ID = 3};
\texttt{R\_L} (resistencia de carga del m\'odulo: ``203'' = 20~k$\Omega$, ``102'' = 1~k$\Omega$) y
\texttt{V\_C} (voltaje real del m\'odulo medido con mult\'imetro). Tras 24~h de precalentamiento,
escribir \texttt{CAL} en el Monitor Serie con el sal\'on ventilado.
\end{completar}

\subsection{Nodo de iluminaci\'on (BH1750)}

\textbf{Conexi\'on.} I$^2$C (SDA al GPIO~21, SCL al GPIO~22), alimentaci\'on a 3.3~V,
direcci\'on 0x23.

\textbf{Funcionamiento.} El BH1750 entrega directamente la iluminancia en lux (1 a 65\,535~lx).
El nodo lee una vez por segundo y env\'ia cada 10~s el promedio, el m\'inimo y el m\'aximo de la
ventana. Como el sensor est\'a en el techo, la iluminancia sobre las mesas se estima con un
factor medido con un lux\'ometro:
\begin{equation}
  E_{\text{mesa}} \approx k\,\bar{E}_{\text{techo}}, \qquad k = 0.482.
\end{equation}
Para detectar el encendido o apagado de las luces se compara el promedio de los \'ultimos 3~s con
el de los 3~s anteriores: el cambio debe superar 50~lx y el 30\,\% del nivel mayor, y ambos bloques
deben ser estables. As\'i se distingue un cambio de luces (inmediato) de una nube (gradual).

\textbf{Variables:} \texttt{lux}, \texttt{lux\_min}, \texttt{lux\_max}, \texttt{lux\_mesa\_est},
\texttt{luces\_cambio}. \textbf{Anomal\'ia} \texttt{ILUMINACION\_BAJA}: $E_{\text{mesa}} < 300$~lx
(se apaga sobre 330~lx). \textbf{Archivo:} \arch{nodo\_luz.ino}.

\begin{completar}
En \arch{nodo\_luz.ino}: \texttt{MAC\_GATEWAY}, \texttt{SSID\_ROUTER}, \texttt{NODO\_ID = 4} y
\texttt{FACTOR\_MESA}. El factor solo es v\'alido para la posici\'on y orientaci\'on en que se midi\'o.
\end{completar}

\subsection{Nodo de conteo de personas (MLX90640)}

\textbf{Conexi\'on.} I$^2$C a 1~MHz (SDA al GPIO~21, SCL al GPIO~22), alimentaci\'on a 3.3~V.
El sensor se instala en el marco superior de la puerta, a 2.3~m, mirando hacia abajo.

\textbf{Geometr\'ia.} Con la cabeza a unos 0.6~m del sensor, el \'area vista a esa altura es
\begin{equation}
  L = 2\,d\,\tan\!\left(\frac{\theta}{2}\right):\quad
  L_{55^\circ} \approx 0.62~\text{m},\quad L_{35^\circ} \approx 0.38~\text{m},
\end{equation}
de modo que cada p\'ixel mide unos 2~cm y una cabeza ocupa alrededor de $9 \times 11$ p\'ixeles.
El eje de 55$^\circ$ se orienta en el sentido de la marcha para ver a cada persona en unos 4 cuadros.

\textbf{Algoritmo.} En cada cuadro de $32\times24$ p\'ixeles: (1) se resta un modelo del fondo
aprendido con la puerta despejada; (2) se marcan los p\'ixeles que superan el fondo en
$\max(1.5~^\circ\text{C},\,4\sigma)$; (3) se agrupan en regiones; (4) cada m\'aximo local de
temperatura se interpreta como una cabeza, lo que permite separar dos personas que pasan juntas;
(5) se sigue cada cabeza entre cuadros y se cuenta un cruce cuando su trayectoria atraviesa dos
l\'ineas virtuales. La ocupaci\'on se acumula como
\begin{equation}
  N(t) = \max\big(0,\; N(t^-) + E - S\big).
\end{equation}

\textbf{Variables:} \texttt{entradas}, \texttt{salidas}, \texttt{ocupacion\_est},
\texttt{personas\_en\_vista}, \texttt{cruces\_dudosos}, \texttt{t\_fondo}, \texttt{t\_max}.
\textbf{Archivo:} \arch{nodo\_termico.ino}.

\begin{completar}
En \arch{nodo\_termico.ino}: \texttt{MAC\_GATEWAY}, \texttt{SSID\_ROUTER}, \texttt{NODO\_ID = 5}.
Con \texttt{DEBUG\_MAPA 1} verificar que la cabeza cruce las l\'ineas de izquierda a derecha y
ajustar \texttt{SENTIDO\_ENTRADA}. Al encender, la puerta debe quedar despejada unos 3~s.
\end{completar}

\subsection{Nodo de la puerta (KY-024)}

\textbf{Conexi\'on.} Salida anal\'ogica AO al GPIO~34 (ADC1), alimentaci\'on a 3.3~V. El sensor
va en el marco y el im\'an en la hoja de la puerta. Opcionalmente, una resistencia de 100~k$\Omega$
entre AO y 3.3~V permite detectar un cable suelto.

\textbf{Funcionamiento.} La salida del KY-024 es invertida: con el im\'an cerca (puerta cerrada)
la lectura es baja. El estado se decide con hist\'eresis (cerrada bajo 1200, abierta sobre 1800,
y entre ambos se mantiene el estado anterior) y un antirrebote de 150~ms. El nodo lee 20 veces por
segundo y env\'ia cada 10~s, y de inmediato cuando la puerta cambia de estado.

\textbf{Variables:} \texttt{puerta\_abierta}, \texttt{aperturas}, \texttt{seg\_abierta},
\texttt{campo\_raw}. \textbf{Anomal\'ia} \texttt{PUERTA\_ABIERTA\_PROLONGADA}: m\'as de 10~min.
\textbf{Archivo:} \arch{nodo\_magnetico.ino}.

\begin{completar}
En \arch{nodo\_magnetico.ino}: \texttt{MAC\_GATEWAY}, \texttt{SSID\_ROUTER}, \texttt{NODO\_ID = 6}.
Con \texttt{DEBUG\_RAW 1} comprobar que la lectura cerrada quede claramente bajo
\texttt{UMBRAL\_CERRADA} y la abierta sobre \texttt{UMBRAL\_ABIERTA}.
\end{completar}

\subsection{Nodo de ambiente (BME280)}

\textbf{Conexi\'on.} I$^2$C (SDA al GPIO~21, SCL al GPIO~22), alimentaci\'on a 3.3~V, direcci\'on
0x76 (SDO a GND). Debe separarse del ESP32 y, sobre todo, del calefactor del MQ-2.

\textbf{Funcionamiento.} El sensor mide en modo forzado cada 10~s (solo cuando se le pide), lo
que minimiza su autocalentamiento. Las anomal\'ias usan hist\'eresis: confort entre 20 y 26~$^\circ$C
(se apaga 0.5~$^\circ$C dentro del rango), humedad alta sobre 70\,\% (se apaga bajo 65\,\%) y cambio
r\'apido si la temperatura var\'ia m\'as de 1.5~$^\circ$C respecto a hace 5~min.

\textbf{Variables:} \texttt{temp\_c}, \texttt{hum\_pct}, \texttt{pres\_hpa}.
\textbf{Archivo:} \arch{nodo\_ambiente.ino}.

\begin{completar}
En \arch{nodo\_ambiente.ino}: \texttt{MAC\_GATEWAY}, \texttt{SSID\_ROUTER}, \texttt{NODO\_ID = 7}.
En Lima la humedad relativa suele estar entre 70 y 85\,\%: tras unos d\'ias de datos conviene
ajustar \texttt{HUM\_ALTA}. Si el monitor indica BMP280, el m\'odulo no mide humedad.
\end{completar}

% =====================================================================
\section{Fase 2: Enlace ESP-NOW y gateway}

\subsection{Obtenci\'on de la MAC del gateway}
\label{sec:mac}

Los nodos necesitan la direcci\'on MAC del gateway para enviarle paquetes. Hay dos formas de
obtenerla:
\begin{enumerate}
  \item Subir cualquier programa al ESP32 v4: la consola de salida de Arduino IDE (no el
        Monitor Serie) muestra la l\'inea \texttt{MAC: xx:xx:xx:xx:xx:xx}.
  \item Subir el programa \arch{obtener\_mac.ino}, que la imprime cada 2~s en el Monitor
        Serie (115\,200 baudios) con el formato listo para copiar.
\end{enumerate}

\begin{lstlisting}[style=cpp, caption={Contenido de \texttt{obtener\_mac.ino}.}]
#include <WiFi.h>

void setup() {
  Serial.begin(115200);
  WiFi.mode(WIFI_STA);   // necesario para que la MAC no salga en ceros
  delay(100);
}

void loop() {
  uint8_t m[6];
  WiFi.macAddress(m);
  Serial.printf("MAC STA: %s\n", WiFi.macAddress().c_str());
  Serial.printf("Para MAC_GATEWAY: {0x%02X, 0x%02X, 0x%02X, 0x%02X, 0x%02X, 0x%02X}\n\n",
                m[0], m[1], m[2], m[3], m[4], m[5]);
  delay(2000);
}
\end{lstlisting}

La MAC obtenida fue \texttt{58:2A:BD:72:03:DC}. Es la MAC en modo estaci\'on (STA), que
tambi\'en es la que ve el punto de acceso WiFi; por eso es la que se entrega a TI.

\subsection{Sincronizaci\'on de canal}

ESP-NOW solo funciona si emisor y receptor est\'an en el mismo canal. Al conectarse al WiFi,
el gateway queda en el canal que elige el router. Por eso cada nodo, al arrancar, escanea la
red indicada en \texttt{SSID\_ROUTER}, lee su canal y se sintoniza en \'el sin conectarse.
Si fallan 10 env\'ios seguidos, vuelve a escanear por si el router cambi\'o de canal. Esta
l\'ogica est\'a en \arch{espnow\_nodo.h}, que no requiere modificaciones.

\begin{nota}[Orden de encendido]
Encender primero el gateway y despu\'es los nodos, para que al escanear la red ya est\'e en su
canal definitivo. La red debe ser de \textbf{2.4~GHz}: el ESP32 no detecta redes de 5~GHz.
\end{nota}

\subsection{C\'odigo del gateway}

El gateway (\arch{gateway.ino} y \arch{paquetes.h}, carpeta \arch{gateway}) realiza:
\begin{enumerate}
  \item Conexi\'on al WiFi y sincronizaci\'on de la hora por NTP (UTC).
  \item Recepci\'on ESP-NOW: el callback solo copia el paquete a una cola de FreeRTOS, porque
        se ejecuta dentro de la tarea del WiFi y no debe bloquearse.
  \item Procesamiento en el \texttt{loop()}: valida versi\'on, tipo y tama\~no; asigna la hora
        real; cuenta paquetes perdidos por nodo a partir de \texttt{seq}.
  \item Conversi\'on a JSON y publicaci\'on en el t\'opico
        \texttt{utec/gemelo/<grupo>/nodo<id>/<tipo>}, con los siete tipos definidos
        en \arch{paquetes.h}.
\end{enumerate}

Las anomal\'ias se env\'ian como texto (\texttt{"CO2\_ALERTA\_VENTILACION|..."}) y como
n\'umero (\texttt{anom\_bits}), porque InfluxDB guarda textos y n\'umeros pero no listas.

\begin{lstlisting}[style=plano, caption={Ejemplo de mensaje publicado por el gateway.}]
utec/gemelo/grupo1/nodo1/co2
{"nodo":1,"mac":"00:70:07:A1:E0:24","seq":1409,"ts":1793552417884,
 "uptime_ms":7050021,"perdidos":0,"tipo":"co2","co2_ppm":571,
 "delta_30s":0,"temp_c":32,"calentando":0,"nivel":"NORMAL",
 "anomalias":"NINGUNA","anom_bits":0}
\end{lstlisting}

\begin{completar}
En \arch{gateway.ino}:
\begin{itemize}
  \item \texttt{WIFI\_SSID} y \texttt{WIFI\_PASS}: red WiFi de 2.4~GHz. Debe ser la misma
        red indicada en \texttt{SSID\_ROUTER} de ambos nodos.
  \item \texttt{GRUPO}: identificador en los t\'opicos. Debe ser el mismo que usa el backend del
        gemelo 2D y, en la nube, el usuario MQTT (\texttt{grupo1}).
  \item \texttt{USAR\_MQTT}: \texttt{0} para validar solo el enlace (salida por Monitor
        Serie); \texttt{1} para publicar en el servidor (Fase~3).
  \item \texttt{MQTT\_HOST}: direcci\'on IPv4 de la laptop (Fase~3, paso~\ref{paso:ip}).
\end{itemize}
\end{completar}

\subsection{Conexi\'on a la red de la universidad}

La red \emph{UTEC Proyectos} usa seguridad WPA2-Personal (clave compartida) y es de doble banda.
El ESP32 solo opera en 2.4~GHz, mientras que una laptop puede conectarse a la misma red en 5~GHz;
ambas quedan en la misma subred. Si el gateway no conecta, el c\'odigo de estado de
\texttt{WiFi.status()} orienta el diagn\'ostico: 1 indica que no encuentra la red (nombre mal
escrito o sin cobertura de 2.4~GHz), 4 que la conexi\'on fue rechazada (clave o MAC no autorizada)
y 6, casi siempre, una clave incorrecta. El ejemplo \emph{WiFiScan} de Arduino IDE confirma si la
red es visible en 2.4~GHz. Si el gateway conecta pero no llega al servidor, la causa probable es el
aislamiento de clientes de la red, que solo puede desactivar TI.

\subsection{Validaci\'on del enlace}

Con \texttt{USAR\_MQTT = 0}, el Monitor Serie del gateway mostr\'o los mensajes de ambos nodos:
audio cada 1~s y CO$_2$ cada 5~s. El estado report\'o \texttt{WiFi=OK}, \texttt{canal=11},
\texttt{NTP=OK} y \texttt{descartados=0}. El contador \texttt{perdidos} se mantuvo en 0 para
el nodo de CO$_2$; el de audio registr\'o 4 p\'erdidas durante el arranque y luego se mantuvo
estable.

% =====================================================================
\section{Fase 3: Servidor local en la laptop}

\subsection{Instalaci\'on de Docker Desktop y WSL~2}

\begin{enumerate}
  \item Descargar Docker Desktop desde \url{https://www.docker.com}. En el instalador,
        elegir \emph{Per-user installation} con backend \textbf{WSL~2} (no \emph{Docker VMM}).
  \item Verificar en el Administrador de tareas $\rightarrow$ Rendimiento $\rightarrow$ CPU
        que la virtualizaci\'on est\'e habilitada.
  \item En PowerShell como administrador, instalar o actualizar WSL y reiniciar:
\begin{lstlisting}[style=bash]
wsl --install        # o:  wsl --update
\end{lstlisting}
  \item Comprobar la versi\'on (debe indicar \texttt{Default Version: 2}):
\begin{lstlisting}[style=bash]
wsl --status
\end{lstlisting}
  \item Abrir Docker Desktop y esperar el indicador verde \emph{Engine running}.
  \item Prueba final:
\begin{lstlisting}[style=bash]
docker --version
docker run hello-world
\end{lstlisting}
\end{enumerate}

\subsection{Estructura del servidor local}

Los archivos se ubican en la carpeta \arch{servidor\_iot}:

\begin{lstlisting}[style=plano]
servidor_iot/
|-- .env
|-- docker-compose.yml
|-- mosquitto/config/mosquitto.conf
|-- telegraf/telegraf.conf
|-- grafana/provisioning/datasources/influxdb.yml
\end{lstlisting}

\begin{table}[H]
\centering
\caption{Archivos del servidor local.}
\begin{tabularx}{\textwidth}{@{}lX@{}}
\toprule
\textbf{Archivo} & \textbf{Funci\'on} \\
\midrule
\arch{docker-compose.yml} & Define los cuatro servicios y su red interna. Expone los puertos
1883 (MQTT), 8086 (InfluxDB) y 3000 (Grafana). \\
\arch{.env} & Contrase\~nas y token de InfluxDB. \\
\arch{mosquitto.conf} & Habilita el listener 1883 con acceso an\'onimo (solo para pruebas
locales). \\
\arch{telegraf.conf} & Se suscribe a los siete tipos (\texttt{utec/gemelo/+/+/co2},
\texttt{audio}, \texttt{gas}, \texttt{luz}, \texttt{termico}, \texttt{magnetico} y \texttt{ambiente}),
interpreta el JSON y escribe en InfluxDB. \\
\arch{influxdb.yml} & Crea autom\'aticamente la fuente de datos InfluxDB en Grafana. \\
\bottomrule
\end{tabularx}
\end{table}

En \arch{telegraf.conf} se define el contrato con el gateway: \texttt{tipo} determina la
tabla (\emph{measurement}) \texttt{co2} o \texttt{audio}; \texttt{nodo} y \texttt{mac} son
etiquetas indexadas; \texttt{ts} es la marca de tiempo (hora del gateway, no la de
llegada); y el grupo se extrae del t\'opico como etiqueta \texttt{grupo}. Los tipos se listan
expl\'icitamente, en lugar de usar el comod\'in \texttt{\#}, porque el t\'opico de estado del gateway
contiene texto plano (\texttt{online}/\texttt{offline}) y no JSON.

\begin{completar}
En \arch{servidor\_iot/.env}, antes del primer arranque:
\begin{itemize}
  \item \texttt{INFLUX\_PASSWORD} y \texttt{GRAFANA\_PASSWORD}: m\'inimo 8 caracteres, sin
        espacios ni comillas.
  \item \texttt{INFLUX\_TOKEN}: texto largo y aleatorio.
\end{itemize}
InfluxDB lee estos valores solo la primera vez. Si se cambian despu\'es, hay que borrar los
vol\'umenes con \texttt{docker compose down -v} (se pierden los datos).
\end{completar}

\subsection{Levantar el servidor}

En una terminal abierta dentro de \arch{servidor\_iot}:
\begin{lstlisting}[style=bash]
docker compose up -d
docker compose ps
\end{lstlisting}
Los cuatro contenedores deben aparecer como \emph{Up}. InfluxDB queda en
\url{http://localhost:8086} y Grafana en \url{http://localhost:3000} (usuario \texttt{admin}).

\subsection{Conectar el gateway al servidor local}

\begin{enumerate}
  \item Conectar la laptop a la misma red WiFi que el gateway.
  \item \label{paso:ip} Obtener la IPv4 de la laptop con \texttt{ipconfig}, en la secci\'on
        \emph{Adaptador de LAN inal\'ambrica Wi-Fi}. No usar las direcciones de los
        adaptadores \emph{vEthernet (WSL)} ni de Docker.
  \item Cambiar el perfil de la red a \textbf{Privada} (Configuraci\'on $\rightarrow$ Red e
        Internet $\rightarrow$ Wi-Fi $\rightarrow$ red $\rightarrow$ Tipo de perfil).
  \item Abrir el puerto 1883 en el firewall (PowerShell como administrador):
\begin{lstlisting}[style=bash]
New-NetFirewallRule -DisplayName "MQTT 1883" -Direction Inbound -Protocol TCP -LocalPort 1883 -Action Allow
\end{lstlisting}
  \item Instalar la librer\'ia \emph{PubSubClient} de Nick O'Leary en Arduino IDE.
  \item En \arch{gateway.ino}: \texttt{USAR\_MQTT = 1} y \texttt{MQTT\_HOST} con la IPv4
        obtenida. Subir el programa; el Monitor Serie debe mostrar \texttt{\# MQTT conectado}.
\end{enumerate}

\subsection{Verificaci\'on de la cadena de datos}

\textbf{Mosquitto.} Desde cualquier terminal:
\begin{lstlisting}[style=bash]
docker exec -it mosquitto mosquitto_sub -t "utec/gemelo/#" -v
\end{lstlisting}
Se observan los mensajes de \texttt{nodo1/co2}, \texttt{nodo2/audio} y
\texttt{gateway/estado}. Este comando es solo un visor: cerrarlo no detiene el servidor.

\textbf{InfluxDB.} En \emph{Data Explorer}, bucket \texttt{gemelo}, se verificaron las
tablas \texttt{co2} y \texttt{audio}, sus campos, la etiqueta \texttt{grupo1} y las MAC de
los nodos. InfluxDB almacena en UTC; la interfaz muestra la hora local (UTC$-5$).

\subsection{Dashboard en Grafana}

El dashboard se importa desde \arch{dashboard\_gemelo\_L418\_final.json}: Grafana $\rightarrow$
Dashboards $\rightarrow$ New $\rightarrow$ Import $\rightarrow$ subir el archivo y elegir la
fuente \emph{InfluxDB}. Se actualiza cada 10~s.

\begin{table}[H]
\centering
\caption{Secciones del dashboard final.}
\begin{tabularx}{\textwidth}{@{}lX@{}}
\toprule
\textbf{Secci\'on} & \textbf{Contenido} \\
\midrule
Estado actual & Ocupaci\'on, puerta, CO$_2$, nivel de gas, $L_{eq}$, temperatura, humedad y luz sobre la mesa \\
Salud de los nodos & Segundos desde el \'ultimo dato y paquetes perdidos de cada nodo \\
Ocupaci\'on y puerta & Ocupaci\'on en el tiempo, entradas y salidas por minuto, franja de la puerta \\
Calidad del aire & CO$_2$, \'indice Rs/R0, GLP equivalente y franja de niveles de CO$_2$ y gas \\
Confort t\'ermico & Temperatura, humedad y presi\'on \\
Ac\'ustica & $L_{eq}$ (promedio energ\'etico), $L_{max}$ y ruido de fondo \\
Iluminaci\'on & Luz estimada sobre la mesa y luz medida en el techo \\
Registro de anomal\'ias & Tabla \'unica de los siete sensores, con filtro por columna \\
Exportar datos crudos & Fila plegada con una tabla por sensor, sin promediar \\
\bottomrule
\end{tabularx}
\end{table}

\textbf{Filtro por grupo.} La variable \emph{Grupo} lista autom\'aticamente los grupos
presentes en la base de datos; cuando otros grupos publiquen, aparecer\'an sin cambiar el
dashboard.

\textbf{Exportaci\'on CSV.} El usuario elige el rango de tiempo, abre la fila
\emph{Exportar datos crudos} y, en el panel deseado, selecciona \emph{Inspect $\rightarrow$
Data $\rightarrow$ Download CSV}. Se exporta desde esa fila porque las gr\'aficas agregan los
datos seg\'un el rango (en 7 d\'ias cada punto representa varios minutos); la fila de
exportaci\'on entrega cada medici\'on original.

\textbf{Promedio energ\'etico en Flux.} La curva de $L_{eq}$ convierte cada valor a
energ\'ia, promedia y vuelve a decibeles:
\begin{lstlisting}[style=plano]
|> map(fn: (r) => ({r with _value: math.pow(x: 10.0, y: r._value / 10.0)}))
|> aggregateWindow(every: v.windowPeriod, fn: mean, createEmpty: false)
|> map(fn: (r) => ({r with _value: 10.0 * math.log10(x: r._value)}))
\end{lstlisting}

\begin{completar}
\arch{dashboard\_gemelo\_L418\_final.json} no requiere cambios. Al importarlo solo se
selecciona la fuente de datos InfluxDB; si ya exist\'ia una versi\'on anterior, se elige
\emph{Import (Overwrite)}.
\end{completar}

% =====================================================================
\section{Fase 4: Migraci\'on a la nube}

\subsection{Cambios necesarios para exponer el servidor a internet}

Los contenedores son los mismos; lo que cambia es la seguridad, porque el servidor queda
accesible desde internet.

\begin{table}[H]
\centering
\caption{Puertos expuestos en la nube.}
\begin{tabularx}{\textwidth}{@{}llX@{}}
\toprule
\textbf{Puerto} & \textbf{Servicio} & \textbf{Protecci\'on} \\
\midrule
80 / 443 & Caddy $\rightarrow$ Grafana & HTTPS con certificado Let's Encrypt obtenido y
renovado autom\'aticamente. \\
8883 & Mosquitto & TLS, usuario y contrase\~na, permisos por grupo. \\
22 & SSH & Solo desde la IP del administrador; acceso con llave privada. \\
\midrule
8086 & InfluxDB & No expuesto. Acceso solo por t\'unel SSH. \\
3000, 1883 & Grafana, MQTT sin cifrar & No expuestos. \\
\bottomrule
\end{tabularx}
\end{table}

\begin{itemize}
  \item \textbf{Grafana p\'ublico de solo lectura.} Se habilita el acceso an\'onimo con rol
        \emph{Viewer}: cualquier persona ve el dashboard, cambia el rango y descarga CSV, pero
        no edita. No se usan los \emph{public dashboards} de Grafana porque no admiten la
        variable \emph{Grupo}.
  \item \textbf{MQTT por grupos.} Cada gateway usa como usuario el nombre de su grupo y, por
        la lista de control de acceso, solo puede publicar en
        \texttt{utec/gemelo/<su grupo>/\#}. Telegraf solo puede leer.
  \item \textbf{TLS con autoridad certificadora propia.} El certificado de MQTT lo firma una
        CA propia v\'alida por 10 a\~nos, grabada en el gateway. As\'i el gateway solo acepta
        este servidor y no necesita reprogramarse cuando se renuevan los certificados
        HTTPS (cada 90~d\'ias).
\end{itemize}

\subsection{Creaci\'on de la m\'aquina virtual en Azure for Students}

Azure for Students otorga 100~USD de cr\'edito anual sin tarjeta de cr\'edito. Los servicios
gratuitos incluyen 750 horas mensuales de los tama\~nos B1s, B2pts~v2 y B2ats~v2, suficiente
para una VM encendida todo el mes (744 horas en un mes de 31 d\'ias).

\begin{enumerate}
  \item Activar la oferta en \url{https://azure.microsoft.com/free/students} con el correo
        de UTEC y verificar en \emph{Suscripciones} que aparezca \emph{Azure for Students}.
  \item En el portal, buscar \textbf{M\'aquinas virtuales} (no \emph{Operator Nexus}) y
        elegir \emph{Crear $\rightarrow$ M\'aquina virtual de Azure}.
\end{enumerate}

\begin{table}[H]
\centering
\caption{Configuraci\'on de la m\'aquina virtual.}
\begin{tabularx}{\textwidth}{@{}lX@{}}
\toprule
\textbf{Campo} & \textbf{Valor} \\
\midrule
Grupo de recursos & \texttt{rg-gemelo} \\
Nombre & \texttt{vm-gemelo} \\
Regi\'on & Una de las permitidas por la suscripci\'on donde B2ats~v2 est\'e disponible
(North Central US o Mexico Central) \\
Imagen & Ubuntu Server 24.04 LTS -- x64 Gen2 \\
Tama\~no & Standard\_B2ats\_v2 (2 vCPU, 1~GiB), apto para servicios gratuitos \\
Autenticaci\'on & Clave p\'ublica SSH, usuario \texttt{azureuser}, par \texttt{vm-gemelo\_key} \\
Disco del SO & 64~GiB, SSD Premium (P6) \\
Puertos de entrada & Solo SSH (22) al crear \\
Apagado autom\'atico & Desactivado \\
\bottomrule
\end{tabularx}
\end{table}

\begin{enumerate}[resume]
  \item En \emph{Revisar y crear}, confirmar que el tama\~no sea B2ats~v2 (el valor por
        defecto, D2s~v3, no es gratuito). Al crear, descargar la llave privada
        \arch{vm-gemelo\_key.pem}; no se puede volver a descargar.
  \item En la VM $\rightarrow$ Informaci\'on general $\rightarrow$ \emph{Nombre DNS},
        asignar una etiqueta, por ejemplo \texttt{gemelo-l418}. La direcci\'on resultante,
        del tipo \texttt{gemelo-l418.<regi\'on>.cloudapp.azure.com}, reemplaza la compra de
        un dominio.
  \item En \emph{Redes}, crear reglas de entrada para los puertos 80, 443 y 8883, y
        restringir el origen de la regla SSH a la IP del administrador.
\end{enumerate}

\begin{nota}[Estado actual]
Durante la creaci\'on se encontraron dos restricciones: (1) la suscripci\'on solo permite las
regiones \texttt{northcentralus}, \texttt{westus}, \texttt{canadacentral}, \texttt{westus2} y
\texttt{mexicocentral}; y (2) una directiva de la Gerencia de TI de UTEC bloquea la creaci\'on
de recursos para proyectos no autorizados. Se solicit\'o a TI, a trav\'es del profesor y de
Mesa de Ayuda, habilitar la creaci\'on de la VM o proporcionar una m\'aquina virtual
institucional. Los pasos siguientes est\'an listos para ejecutarse cuando se obtenga el
acceso.
\end{nota}

\subsection{Preparaci\'on de la m\'aquina virtual}

\textbf{Conexi\'on por SSH desde la laptop.} Mover la llave a
\arch{C:\textbackslash Users\textbackslash <usuario>\textbackslash .ssh\textbackslash} y
restringir sus permisos (SSH rechaza llaves con permisos abiertos):
\begin{lstlisting}[style=bash]
$key = "C:\Users\Usuario\.ssh\vm-gemelo_key.pem"
icacls $key /inheritance:r
icacls $key /grant:r "$($env:USERNAME):(R)"
ssh -i $key azureuser@<DOMINIO>
\end{lstlisting}

\textbf{Instalaci\'on de Docker y memoria swap} (dentro de la VM):
\begin{lstlisting}[style=bash]
sudo apt update && sudo apt upgrade -y
curl -fsSL https://get.docker.com | sudo sh
sudo usermod -aG docker azureuser

sudo fallocate -l 2G /swapfile
sudo chmod 600 /swapfile
sudo mkswap /swapfile
sudo swapon /swapfile
echo '/swapfile none swap sw 0 0' | sudo tee -a /etc/fstab
\end{lstlisting}
Tras volver a conectarse, verificar con \texttt{docker run hello-world} y \texttt{free -h}.
La swap de 2~GB complementa el 1~GiB de RAM del tama\~no B2ats~v2; en la laptop, el stack
completo us\'o unos 146~MB.

\subsection{Servidor de producci\'on}

Los archivos se encuentran en la carpeta \arch{servidor\_iot\_prod}:

\begin{lstlisting}[style=plano]
servidor_iot_prod/
|-- .env
|-- docker-compose.yml
|-- preparar.sh
|-- caddy/Caddyfile
|-- mosquitto/config/mosquitto.conf
|-- mosquitto/config/acl
|-- telegraf/telegraf.conf
|-- grafana/provisioning/datasources/influxdb.yml
|-- grafana/provisioning/dashboards/dashboards.yml
|-- grafana/dashboards/dashboard_gemelo_iot.json
\end{lstlisting}

\begin{table}[H]
\centering
\caption{Archivos del servidor de producci\'on.}
\begin{tabularx}{\textwidth}{@{}lX@{}}
\toprule
\textbf{Archivo} & \textbf{Funci\'on} \\
\midrule
\arch{docker-compose.yml} & Cinco servicios (se agrega Caddy). L\'imite de tama\~no de logs
para no llenar el disco. InfluxDB publicado solo en \texttt{127.0.0.1}. \\
\arch{Caddyfile} & Proxy HTTPS hacia Grafana con certificado autom\'atico. \\
\arch{mosquitto.conf} & Sin acceso an\'onimo. Listener interno 1883 (Telegraf) y externo
8883 con TLS (gateways). \\
\arch{acl} & Telegraf lee todo; cada gateway solo publica en los t\'opicos de su grupo. \\
\arch{telegraf.conf} & Igual al local, con usuario y contrase\~na MQTT. \\
\arch{dashboards.yml} y \arch{dashboard\_gemelo\_iot.json} & Cargan el dashboard
autom\'aticamente como p\'agina de inicio, con actualizaci\'on cada 10~s. \\
\arch{preparar.sh} & Genera la CA y el certificado TLS del broker, crea los usuarios MQTT y
ajusta permisos. Se ejecuta una vez. \\
\bottomrule
\end{tabularx}
\end{table}

\begin{completar}
En \arch{servidor\_iot\_prod/.env} (contrase\~nas solo con letras y n\'umeros):
\begin{itemize}
  \item \texttt{DOMINIO}: nombre DNS de la VM (por ejemplo
        \texttt{gemelo-l418.northcentralus.cloudapp.azure.com}).
  \item \texttt{ACME\_EMAIL}: correo para avisos de Let's Encrypt.
  \item \texttt{INFLUX\_PASSWORD}, \texttt{INFLUX\_TOKEN}, \texttt{GRAFANA\_ADMIN\_PASSWORD}.
  \item \texttt{MQTT\_TELEGRAF\_PASS} y \texttt{MQTT\_GRUPO1\_PASS}.
\end{itemize}
Para integrar otro grupo, agregar en \arch{preparar.sh} una l\'inea
\texttt{mosquitto\_passwd -b ... grupoN <clave>} y volver a ejecutarlo.
\end{completar}

\textbf{Despliegue.} Copiar la carpeta a la VM desde PowerShell y levantar el servidor:
\begin{lstlisting}[style=bash]
scp -i C:\Users\Usuario\.ssh\vm-gemelo_key.pem -r servidor_iot_prod azureuser@<DOMINIO>:~
\end{lstlisting}
\begin{lstlisting}[style=bash]
cd servidor_iot_prod
bash preparar.sh       # al final imprime el certificado ca.crt
docker compose up -d
\end{lstlisting}

\subsection{Gateway para la nube}

El archivo \arch{gateway\_nube.ino} (carpeta \arch{gateway\_nube}, junto con
\arch{paquetes.h}) es el gateway final con \texttt{USAR\_TLS = 1}: decodifica los siete tipos,
usa \texttt{WiFiClientSecure}, el puerto 8883 y usuario y contrase\~na, y espera a tener la hora
NTP antes de conectarse, porque sin hora real no puede validar la vigencia del certificado. El
usuario MQTT es el mismo valor de \texttt{GRUPO}, como exige la lista de control de acceso.

\begin{completar}
En \arch{gateway\_nube.ino}:
\begin{itemize}
  \item \texttt{WIFI\_SSID} y \texttt{WIFI\_PASS}: red WiFi de la universidad autorizada por TI
        para el gateway (MAC \texttt{58:2A:BD:72:03:DC}). Los nodos deben tener el mismo valor
        en \texttt{SSID\_ROUTER}.
  \item \texttt{MQTT\_HOST}: el \texttt{DOMINIO} del \arch{.env}.
  \item \texttt{GRUPO}: \texttt{grupo1}, igual al usuario creado en \arch{preparar.sh}
        (\texttt{MQTT\_USER} toma este valor autom\'aticamente).
  \item \texttt{MQTT\_PASS}: el valor de \texttt{MQTT\_GRUPO1\_PASS}.
  \item \texttt{CA\_CERT}: el contenido completo de \arch{ca/ca.crt} impreso por
        \arch{preparar.sh}, desde \texttt{-----BEGIN CERTIFICATE-----} hasta
        \texttt{-----END CERTIFICATE-----}.
\end{itemize}
\end{completar}

\subsection{Verificaci\'on en la nube}

\begin{enumerate}
  \item Abrir \texttt{https://<DOMINIO>} desde un celular con datos m\'oviles: el dashboard
        debe cargarse sin iniciar sesi\'on.
  \item Revisar que el Monitor Serie del gateway muestre \texttt{\# MQTT conectado}.
  \item Para administrar InfluxDB, abrir un t\'unel SSH y entrar a
        \url{http://localhost:8086} en la laptop:
\begin{lstlisting}[style=bash]
ssh -i C:\Users\Usuario\.ssh\vm-gemelo_key.pem -L 8086:localhost:8086 azureuser@<DOMINIO>
\end{lstlisting}
\end{enumerate}

\subsection{Control de costos}

En \emph{Cost Management $\rightarrow$ Presupuestos} se cre\'o el presupuesto
\texttt{presupuesto-gemelo}: 10~USD mensuales, vigente hasta octubre de 2027, con alertas al
50\,\% y 80\,\% del gasto real y al 100\,\% del gasto previsto. El presupuesto solo notifica;
si el cr\'edito se agotara, Azure deshabilitar\'ia los recursos sin cargos a tarjeta.

% =====================================================================
\section{Integraci\'on con el gemelo digital 2D}

El gemelo digital 2D (repositorio Web-PI3) es una aplicaci\'on web que muestra el plano del aula
con cada sensor coloreado seg\'un su estado. Su backend (FastAPI) \textbf{solo lee} de InfluxDB:
cada 2~s consulta la \'ultima lectura de cada sensor, calcula su estado a partir del \texttt{nivel}
y las \texttt{anomalias} que env\'ian los nodos, y env\'ia los cambios a los navegadores mediante
\emph{Server-Sent Events} (SSE). Con muchos navegadores abiertos, InfluxDB recibe igualmente una
sola consulta por ciclo.

\begin{table}[H]
\centering
\caption{Requisitos de la arquitectura IoT para alimentar el gemelo 2D.}
\begin{tabularx}{\textwidth}{@{}lX@{}}
\toprule
\textbf{Requisito} & \textbf{Detalle} \\
\midrule
Acceso a InfluxDB & \texttt{http://localhost:8086} si el backend corre en el mismo equipo, o
\texttt{http://influxdb:8086} si se agrega como servicio al \arch{docker-compose.yml} \\
Token de solo lectura & Creado en InfluxDB con permiso \emph{Read} sobre el bucket \texttt{gemelo};
nunca el token de administrador \\
Mismo \texttt{GRUPO} & El backend filtra por la etiqueta \texttt{grupo}; debe coincidir con el
\texttt{GRUPO} del gateway \\
Cat\'alogo completo & El archivo \arch{sensores.yaml} debe declarar los siete sensores con los
nombres exactos de tablas y campos de InfluxDB \\
Demora & Con \texttt{flush\_interval = "1s"} en Telegraf, la demora sensor--plano queda en unos 3~s \\
\bottomrule
\end{tabularx}
\end{table}

% =====================================================================
\section{Resultados de la validaci\'on local}

\begin{table}[H]
\centering
\caption{Resultados observados con el servidor en la laptop.}
\begin{tabularx}{\textwidth}{@{}lX@{}}
\toprule
\textbf{Prueba} & \textbf{Resultado} \\
\midrule
Enlace ESP-NOW & Nodo CO$_2$: 0 paquetes perdidos. Nodo audio: 4 p\'erdidas al arranque y
0 en una sesi\'on posterior de aproximadamente 2~h (m\'as de 6800 paquetes). \\
Sincronizaci\'on & NTP correcto; canal WiFi 11; 0 paquetes descartados en el gateway. \\
Latencia & El indicador de salud mostr\'o entre 2 y 10~s desde el \'ultimo dato. \\
Recursos del servidor & Aproximadamente 146~MB de RAM y 0.5\,\% de CPU para los cuatro
contenedores. \\
CO$_2$ & Entre 560 y 620~ppm, estado Normal. Se observ\'o la respuesta del sensor ante
exhalaci\'on. \\
Audio & Niveles entre 82 y 87~dB con ruido de fondo estimado cercano a 83~dB. El valor
sugiere una fuente de ruido cercana al micr\'ofono o la necesidad de calibrar
\texttt{OFFSET\_SPL\_DB} con un son\'ometro. \\
Exportaci\'on & Descarga CSV verificada; la representaci\'on se adapta al rango de tiempo. \\
Nodos 3 a 7 & Verificados en simulaci\'on: conteo de personas (cruces, fusiones, retrocesos y
falsos positivos), puerta (rebotes, zona intermedia, cable suelto), ambiente, luz (luces frente a
nubes) y gas (niveles con hist\'eresis y persistencia de $R_0$). La prueba integrada de los siete
nodos en el sal\'on se programa para el 17 de noviembre de 2026. \\
\bottomrule
\end{tabularx}
\end{table}

% =====================================================================
\section{Problemas encontrados y soluciones}

\begin{longtable}{@{}p{0.42\textwidth}p{0.53\textwidth}@{}}
\caption{Problemas y soluciones.}\\
\toprule
\textbf{Problema} & \textbf{Soluci\'on} \\
\midrule
\endfirsthead
\toprule
\textbf{Problema} & \textbf{Soluci\'on} \\
\midrule
\endhead
ESP32 no detecta redes de 5~GHz. & Usar una red o hotspot configurado en 2.4~GHz. \\
ESP-NOW falla si emisor y receptor est\'an en canales distintos. & Los nodos escanean la red
del gateway para obtener su canal; re-escaneo tras 10 fallos. \\
El gateway no conecta a MQTT en la laptop. & Usar la IPv4 del adaptador Wi-Fi (no
vEthernet), perfil de red privado y regla de firewall para el puerto 1883. \\
Promediar decibeles directamente da resultados incorrectos. & Promedio energ\'etico de
$L_{eq}$ en la consulta Flux. \\
Error en el panel \emph{Nivel sonoro} (funci\'on \texttt{math.pow10}). & Reemplazo por
\texttt{math.pow(x: 10.0, y: ...)}. \\
Las gr\'aficas exportan valores agregados. & Fila separada de exportaci\'on con datos
crudos. \\
Formulario equivocado en Azure (\emph{Operator Nexus}). & Usar el servicio
\emph{M\'aquinas virtuales} est\'andar. \\
B2ats~v2 no disponible en Brazil South
(\texttt{NotAvailableForSubscription}). & Elegir una regi\'on permitida donde el tama\~no
est\'e disponible. \\
Directiva de regiones permitidas de la suscripci\'on. & Consultar la lista en
\emph{Directiva $\rightarrow$ Tareas (asignaciones) $\rightarrow$ Par\'ametros}. \\
Directiva de TI de UTEC que bloquea la creaci\'on de recursos. & Solicitud a TI y Mesa de
Ayuda a trav\'es del profesor. \\
Pines del ADC2 inutilizables con la radio activa. & Todas las entradas anal\'ogicas usan pines
del ADC1 (GPIO 32 a 39). \\
Salidas de 5~V del MQ-2 frente a un ADC de 3.3~V. & Divisores de tensi\'on de 10/15~k$\Omega$. \\
C\'odigos de nodos con formatos de paquete distintos. & Un \'unico contrato (\arch{paquetes.h}) y
una \'unica tabla de \texttt{NODO\_ID} para todo el sal\'on. \\
Gateway sin conexi\'on a la red de la universidad. & Diagn\'ostico con \emph{WiFiScan} y el c\'odigo
de \texttt{WiFi.status()}; red WPA2-Personal en 2.4~GHz. \\
\bottomrule
\end{longtable}

% =====================================================================
\section{Trabajo pendiente}

\begin{enumerate}
  \item Prueba integrada de los siete nodos en el sal\'on L418 (17 de noviembre de 2026).
  \item Confirmar con TI la comunicaci\'on entre dispositivos en \emph{UTEC Proyectos} y la
        habilitaci\'on de la m\'aquina virtual en la nube.
  \item Calibrar el micr\'ofono con un son\'ometro y, como mejora, aplicar ponderaci\'on~A.
  \item Calibrar $R_0$ del MQ-2 tras 24~h de precalentamiento y verificar $R_L$ y $V_C$.
  \item Completar el cat\'alogo del gemelo 2D con los siete sensores y validar la lectura en tiempo real.
  \item Validar el conteo de personas con conteos manuales en cambios de clase.
\end{enumerate}

% =====================================================================
\section*{Referencias}
\addcontentsline{toc}{section}{Referencias}
\setlength{\parindent}{-1.2em}\setlength{\leftskip}{1.2em}

Adafruit. (s.\ f.). \emph{Adafruit MLX90640 IR Thermal Camera Breakout - 55 Degree: Datasheet} [Product ID: 4407]. Adafruit Industries. Recuperado el 10 de noviembre de 2026, de \url{https://www.adafruit.com/product/4407}

AG Electr\'onica. (2019). \emph{Sensor de presi\'on atmosf\'erica GY-BME280} (No. Parte: OKY3259-2). AG Electr\'onica S.A.P.I. de C.V. Recuperado el 10 de noviembre de 2026, de \url{https://agelectronica.com/detalle?busca=OKY3259-2}

Handson Technology. (s.\ f.). \emph{BH1750 Ambient Light Sensor Module: Data Specs} (SKU: SSR1054). Handson Technology. Recuperado el 10 de noviembre de 2026, de \url{https://handsontec.com/index.php/product/bh1750-ambient-light-sensor-module/}

Joy-IT. (2021). \emph{KY-024 Linear magnetic Hall Sensor: Manual} (SensorKit X40). Joy-IT / Simac Electronics B.V. Recuperado el 10 de noviembre de 2026, de \url{https://joy-it.net/en/products/SEN-KY024LM}

Naylamp Mechatronics. (s.\ f.). \emph{Micr\'ofono I2S INMP441}. Naylamp Mechatronics. Recuperado el 10 de noviembre de 2026, de \url{https://naylampmechatronics.com/audio-sonido/1037-microfono-i2s-inmp441.html}

Technologies Private Limited. (s.\ f.). \emph{ADIY MQ-2 Smoke Sensor Module: Datasheet}. ADIY. Recuperado el 10 de noviembre de 2026, de \url{https://adiy.in/shop/mq2-smoke-sensor-module/}

Zhengzhou Winsen Electronics Technology Co., Ltd. (2023). \emph{Infrared CO2 Sensor (Model: MH-Z1911A): User's Manual} (Version 1.1). Winsen. Recuperado el 10 de noviembre de 2026, de \url{https://www.winsen-sensor.com/product/mh-z1911a.html?searchid=5618}

\setlength{\parindent}{0pt}\setlength{\leftskip}{0pt}

% =====================================================================
\appendix
\section{Repositorio y archivos}

Repositorio del proyecto: \textbf{https://github.com/ARomuloChuquillanqui/gemelo-digital-L418.git}\\
Gemelo digital 2D: \url{https://github.com/MaintzTachyon/Web-PI3}

En el repositorio, el firmware est\'a en \arch{firmware/}, el servidor en \arch{servidor/}, los
dashboards en \arch{grafana/} y este informe en \arch{docs/}.

\begin{longtable}{@{}>{\raggedright\arraybackslash}p{0.30\textwidth}>{\raggedright\arraybackslash}p{0.24\textwidth}>{\raggedright\arraybackslash}p{0.40\textwidth}@{}}
\caption{Archivos del proyecto y datos a completar.}\\
\toprule
\textbf{Archivo} & \textbf{Carpeta} & \textbf{Qu\'e completar} \\
\midrule
\endfirsthead
\toprule
\textbf{Archivo} & \textbf{Carpeta} & \textbf{Qu\'e completar} \\
\midrule
\endhead
\arch{paquetes.h} & todas las de ESP32 & Nada. Id\'entico en todas las copias. \\
\arch{espnow\_nodo.h} & carpetas de los nodos & Nada. \\
\arch{nodo\_co2.ino} & \arch{nodo\_co2} & \texttt{MAC\_GATEWAY}, \texttt{SSID\_ROUTER},
\texttt{NODO\_ID = 1}. \\
\arch{nodo\_audio.ino} & \arch{nodo\_audio} & \texttt{MAC\_GATEWAY}, \texttt{SSID\_ROUTER},
\texttt{NODO\_ID = 2}, \texttt{OFFSET\_SPL\_DB} (calibraci\'on). \\
\arch{nodo\_gas.ino} & \arch{nodo\_gas} & \texttt{MAC\_GATEWAY}, \texttt{SSID\_ROUTER},
\texttt{R\_L}, \texttt{V\_C}; calibrar $R_0$ con \texttt{CAL}. \\
\arch{nodo\_luz.ino} & \arch{nodo\_luz} & \texttt{MAC\_GATEWAY}, \texttt{SSID\_ROUTER},
\texttt{FACTOR\_MESA}. \\
\arch{nodo\_termico.ino} & \arch{nodo\_termico} & \texttt{MAC\_GATEWAY}, \texttt{SSID\_ROUTER},
\texttt{SENTIDO\_ENTRADA}. \\
\arch{nodo\_magnetico.ino} & \arch{nodo\_magnetico} & \texttt{MAC\_GATEWAY},
\texttt{SSID\_ROUTER}, umbrales de puerta. \\
\arch{nodo\_ambiente.ino} & \arch{nodo\_ambiente} & \texttt{MAC\_GATEWAY}, \texttt{SSID\_ROUTER},
\texttt{HUM\_ALTA} (opcional). \\
\arch{obtener\_mac.ino} & \arch{obtener\_mac} & Nada. \\
\arch{gateway.ino} & \arch{gateway} & \texttt{WIFI\_SSID}, \texttt{WIFI\_PASS},
\texttt{GRUPO}, \texttt{USAR\_MQTT}, \texttt{MQTT\_HOST} (IPv4 de la laptop). \\
\arch{.env} & \arch{servidor\_iot} & \texttt{INFLUX\_PASSWORD}, \texttt{INFLUX\_TOKEN},
\texttt{GRAFANA\_PASSWORD}. \\
\arch{docker-compose.yml}, \arch{mosquitto.conf}, \arch{telegraf.conf}, \arch{influxdb.yml}
& \arch{servidor\_iot} & Nada. \\
\arch{dashboard\_gemelo\_\allowbreak L418\_final.json} & \arch{grafana} & Nada; elegir la fuente InfluxDB al
importar. \\
\arch{.env} & \arch{servidor\_iot\_prod} & \texttt{DOMINIO}, \texttt{ACME\_EMAIL},
contrase\~nas de InfluxDB, Grafana y usuarios MQTT. \\
\arch{preparar.sh} & \arch{servidor\_iot\_prod} & Nada; agregar usuarios de otros grupos
cuando se integren. \\
\arch{docker-compose.yml}, \arch{Caddyfile}, \arch{mosquitto.conf}, \arch{acl},
\arch{telegraf.conf}, \arch{influxdb.yml}, \arch{dashboards.yml},
\arch{dashboard\_gemelo\_iot.json} & \arch{servidor\_iot\_prod} & Nada (se carga autom\'aticamente). \\
\arch{gateway\_nube.ino} & \arch{gateway\_nube} & \texttt{WIFI\_PASS}, \texttt{GRUPO},
\texttt{MQTT\_HOST}, \texttt{MQTT\_PASS}, \texttt{CA\_CERT}. \\
\bottomrule
\end{longtable}

\begin{nota}[Seguridad del repositorio]
No subir al repositorio los archivos \arch{.env} con contrase\~nas reales, la llave
\arch{vm-gemelo\_key.pem}, la carpeta \arch{ca/} ni los certificados generados. Publicar en
su lugar versiones con valores de ejemplo.
\end{nota}

\section{Resumen de comandos}

\begin{table}[H]
\centering
\caption{Comandos frecuentes.}
\begin{tabularx}{\textwidth}{@{}>{\ttfamily\small}p{0.55\textwidth}X@{}}
\toprule
\normalfont\textbf{Comando} & \textbf{Uso} \\
\midrule
docker compose up -d & Levantar el servidor (en su carpeta) \\
docker compose ps & Estado de los contenedores \\
docker compose down & Detener sin borrar datos \\
docker compose down -v & Detener y borrar todos los datos \\
docker logs telegraf -{}-tail 20 & Revisar errores de ingesta \\
docker exec -it mosquitto mosquitto\_sub -t "utec/gemelo/\#" -v & Ver mensajes MQTT en vivo
(local) \\
ipconfig & IPv4 de la laptop \\
wsl -{}-shutdown & Reiniciar WSL si Docker no arranca \\
\bottomrule
\end{tabularx}
\end{table}

\end{document}
